\section{The Mass Hierarchy} \label{sec:hierarchy}

One striking feature of the Standard Model remains, the enormous spread of fermion masses, from the electron to the top
quark across five orders of magnitude. We show that this spread is not a set of unrelated small numbers but a single
geometric ladder, a discrete version of the warped extra dimension of Randall and Sundrum \cite{randall1999}, with the
warp factor fixed by the substrate itself.

\subsection{The curved interior and the flat boundary}

The substrate of Section~\ref{sec:law} is hyperbolic, and a hyperbolic space has a sharp two-part structure, a curved
interior, the bulk, and a boundary at infinity, the cusp. The bulk grows exponentially, each shell larger than the last
by a fixed factor. The boundary is flat in the large-scale sense, a flat lattice carrying the field that becomes
ordinary physics.

The physics we observe lives on the flat boundary. The curved interior is not a second place we visit. It is the source
of the warping, the puppet that guides the boundary fields from behind. A field on the boundary has a profile reaching
into the bulk, and because the bulk grows geometrically, that profile is exponentially weighted by depth. A particle
whose profile sits shallow, near the boundary, overlaps strongly with the boundary and is heavy. A particle whose
profile reaches deep into the bulk overlaps weakly and is light. The mass is set by the overlap, and the overlap is set
by the depth, exponentially.

\subsection{The warp factor is the growth rate}

In the continuous Randall-Sundrum picture the warp factor is a free parameter chosen to fit the hierarchy. Here it is
not free. It is the growth rate of the substrate, the asymptotic ratio of successive shell sizes of the $\{3,4,3,4\}$
honeycomb, which we measure by streaming the shell counts outward to large radius. The ratio converges to
\begin{equation}
  \lambda \approx 18.2787 ,
\end{equation}
a pure number fixed by the geometry, the same geometry forced by the octonions. The mass of a fermion localized at
depth $d$ then scales as
\begin{equation}
  m(d) \ \sim \ \lambda^{-c\,d},
\end{equation}
a geometric ladder with common ratio set by $\lambda$. There is also a floor. Normalizability of the boundary profile
requires the overlap not to fall faster than the square root of the growth, so the shallowest, heaviest scale is set by
$\lambda^{1/2} \approx 4.27$, and the ladder of masses descends from there by integer steps of depth.

\begin{result}[The hierarchy is geometric]
The fermion mass hierarchy is a discrete warped-cusp effect. Masses scale as a fixed power $\lambda^{-cd}$ of the
substrate growth rate $\lambda \approx 18.28$, where $d$ is the localization depth of the fermion profile into the
hyperbolic bulk and the heaviest scale is floored at $\lambda^{1/2}$. The large spread of masses is the geometric
spacing of a single ladder, not a set of independent small parameters.
\end{result}

\subsection{What this does and does not fix}

This fixes the \emph{shape} of the hierarchy, that masses fall in a geometric ladder with a ratio fixed by the
geometry, and the scale of the steps. It does not fix which fermion sits at which depth. The assignment of depths, the
integers $d$ for the electron, the muon, the up quark, and so on, is not derived here, just as the individual Yukawa
couplings are not derived in the Standard Model. What is new is that the couplings are no longer arbitrary real
numbers. They are powers of one fixed ratio, so the hierarchy collapses from many independent small numbers to one
geometric scale plus a list of integer depths. We verify the ladder structure and the value of $\lambda$ by direct
shell counting on the substrate, reported in the appendix.

\subsection{Neutrinos and mixing}

Two further features sit naturally on the same ladder, and we record them as consistency checks rather than forced
derivations, since they depend on the depth assignments that the framework leaves open.

The neutrinos are the $N = 0$ color-singlet states, the shallowest in charge but, in the observed spectrum, the
lightest in mass. The geometric ladder accommodates this through the seesaw \cite{randall1999} reading natural to a
warped bulk. A right-handed partner localized very deep in the interior has a large bulk mass, and the light physical
neutrino mass is suppressed by the square of the small overlap, $m_\nu \sim m_D^2 / M$, so the same exponential depth
weighting that spreads the charged fermions pushes the neutrinos far below them. We check that a depth ladder
consistent with the charged-fermion steps reproduces a neutrino scale many orders below the electron, in line with the
smallness of neutrino mass, without a new mechanism.

Flavor mixing fits the same picture. Because each generation is a triality copy of the same eight states, the mixing
between generations is governed by the overlaps of profiles localized at different depths. Profiles at nearby depths
overlap strongly and mix substantially, profiles far apart in depth mix weakly. This gives a mixing pattern that is
near-diagonal with off-diagonal entries falling geometrically in depth separation, qualitatively the hierarchical quark
mixing and the larger lepton mixing seen in nature. We verify the qualitative pattern, near-diagonal with
geometrically suppressed off-diagonals, on the substrate. We do not claim the measured angles, which again require the
depth assignments. The point is only that the mixing is the geometry of overlapping ladders, not a separate
structure.
